\section*{Hardness for deterministic 3-round protocols (author Section 4)}
Building on the proof ideas presented in Section 3, in this section, we prove NP-hardness of approximating the communication complexity of deterministic 3-round protocols. A key building block is to use the notion of universal set.
\setcounter{theorem}{14}
\begin{definition}[Universal set]
Let $r,c,k\in\mathbb N$, and let $M\in\{0,1\}^{r\times c}$ be a matrix whose columns are $M^{(1)},\ldots,M^{(c)}$. We say that $M$ is $(c,k)$-universal if, for every set $\{y_1,\ldots,y_k\}\subseteq[c]$ of $k$ columns of $M$, the matrix
\[\begin{bmatrix}M^{(y_1)}&\cdots&M^{(y_k)}\end{bmatrix}\]
has $2^k$ distinct rows. The set of all the rows of $M$ is called a $(c,k)$-universal set.
\end{definition}
We use the explicit construction of a universal set due to \cite{ref49}.
\begin{lemma}[Naor, Schulman and Srinivasan \cite{ref51}]
There exists a deterministic algorithm that, given $c$ and $k\in\mathbb N$, outputs a $(c,k)$-universal matrix $M\in\{0,1\}^{r\times c}$ such that $r=2^{k+O(\log^2 k)}\cdot\log c$ in time a polynomial in $c$ and $2^k$.
\end{lemma}
Now we state the main result of this section.
\begin{theorem}\label{thm:boolean-construction}
Let $\epsilon>0$ be an arbitrary constant. Given an undirected graph $G=([n],E)$ with $n$ vertices and $|E|=m>0$ edges, one may construct in deterministic polynomial time a function $f_G:[a]\times[b]\to\{0,1\}$ and a number $\ell\in\mathbb N$, with $a,b,\ell=O(n^8)$, such that
\[\chi(G)\cdot\ell\le L_3^B(f_G)\le\chi(G)\cdot\ell^{1+\epsilon}\]
and
\[\log\chi(G)+\log\ell\le C_3^B(f_G)\le\log\chi(G)+(1+\epsilon)\cdot\log\ell.\]
\end{theorem}
The idea of the proof is to build upon the reduction in the proof of Theorem~12, by replacing each column with entries from $[\ell]$ with a block of columns that have entries from $\{0,1\}$. The difficulty in making this work is that a partition of the columns might not respect our blocks and could place columns from the same block into different parts. We solve this by thinking as follows. Either the partition is large, meaning it has many parts, so the protocol also has many leaves, which proves our lower bound, or otherwise for any block $C_v$ of columns the part $P_i$ which has most columns from $C_v$ has many columns from $C_v$; we may then act as if ``$C_v$ was placed in $P_i$''.
\begin{proof}
Let $t,k,c\in\mathbb N$ be parameters chosen later. Let $A\in\{0,1\}^{r\times c}$ and $M\in\{0,1\}^{s\times c}$ be the $(c,t)$-universal and $(c,k)$-universal matrices, respectively, that are constructed by the polynomial-time deterministic algorithm of Lemma~16; then we have $r=2^{(1+o(1))\cdot t}\cdot\log c$ and $s=2^{(1+o(1))\cdot k}\cdot\log c$.
Let $x_1,\ldots,x_s$ be the rows of $M$. We then let $B$ and $C$ denote the $s^2\times c$ matrices
\[
B=\begin{bmatrix}x_1\\\vdots\\x_1\\\vdots\\x_s\\\vdots\\x_s\end{bmatrix},\qquad
C=\begin{bmatrix}x_1\\\vdots\\x_s\\\vdots\\x_1\\\vdots\\x_s\end{bmatrix},
\]
where a row $x_i$ appears $s$ times for each $i\in[s]$. Given an edge $\{v,w\}\in[n]\times[n]$, with $v<w$, we define the $s^2\times c\cdot n$ matrix:
\[M_{\{v,w\}}=\begin{bmatrix}0&\cdots&0&B&0&\cdots&0&C&0&\cdots&0\end{bmatrix},\]
where $B$ appears in the $v$-th block of $c$ columns and $C$ in the $w$-th block of $c$ columns. Finally, let $E=\{e_1,\ldots,e_m\}$ denote the edges of $G$; then we define the $(r+m\cdot s^2)\times c\cdot n$ communication matrix $f_G$ so that
\[f_G=\begin{bmatrix}A&\cdots&A\\\multicolumn{3}{c}{M_{e_1}}\\\multicolumn{3}{c}{\vdots}\\\multicolumn{3}{c}{M_{e_m}}\end{bmatrix}\in\{0,1\}^{(r+m\cdot s^2)\times c\cdot n}.\]
Let $\ell:=2^t$. Let $\mathcal P$ be a partition of the columns of $f_G$ that minimizes $\operatorname{Cost}_{f_G}(\mathcal P)$. We now show that, for a suitable choice of $t,k$ and $c$,
\[\chi(G)\cdot\ell\le\operatorname{Cost}_{f_G}(\mathcal P)\le\chi(G)\cdot\ell^{1+o(1)}.\]
This will complete the proof, because Proposition~11 shows that, for a binary matrix $f_G$, $\operatorname{Cost}_{f_G}(\mathcal P)$ is a 2-factor approximation of $L_3^B(f_G)$ and $\log\operatorname{Cost}_{f_G}(\mathcal P)$ is an approximation of $C_3^B(f_G)$ up to an additive $O(1)$ term. We will choose the parameters $t,k$ and $c$ so that the following conditions are satisfied.
\begin{enumerate}
\item Condition 1. $r+ms\le\ell^{1+o(1)}$.
\item Condition 2. $c\ge n\ell\cdot\max\{t,k\}$.
\item Condition 3. $2^{2k}\ge n\ell$.
\end{enumerate}
Given that Condition 1 is satisfied, the complexity upper-bounds are easy to see. Let $\mathcal P$ be a coloring of $G$ that partitions the vertex set into $\chi(G)$ parts. Since no class contains an edge, the $C$ and $B$ sub-matrices of each $M_{e_i}$ sub-matrix will always be placed in different parts, and we thus have
\[\operatorname{Cost}(\mathcal P)\le\chi(G)\cdot(r+ms)\le\chi(G)\cdot\ell^{1+o(1)}.\]
For the other direction, we will lower-bound $\operatorname{Cost}(\mathcal P)$ for every partition $\mathcal P$. To begin, if the number of parts is $|I|\ge n\ell$, then we must conclude that $\operatorname{Cost}(\mathcal P)\ge n\ell\ge\chi(G)\cdot\ell$, as desired.
Otherwise, $|I|\le n\ell$. Then for every block of columns $v\in[n]$ there must exist a part $P_{i(v)}$ which contains at least $c/|I|\ge c/n\ell$ columns from the $v$-th block. Observe that $c/n\ell\ge\max\{t,k\}$ by Condition 2.
Now, either the mapping $v\mapsto i(v)$ is a valid coloring of $G$ or not. First, suppose that the mapping $v\mapsto i(v)$ is not a valid coloring of $G$, meaning that there exists an edge $\{v,w\}\in E$ such that $i=i(v)=i(w)$. Then the $v$-th column block and the $w$-th column block each have at least $\max\{t,k\}$ columns in the same part $P_i$. This will mean that the $C$ and $B$ sub-matrices of the $M_{\{v,w\}}$ sub-matrix each have at least $k$ columns in $P_i$. But then, since $M$ is $(c,k)$-universal, it follows from Condition 3 that $\operatorname{Cost}(\mathcal P)\ge\operatorname{Cost}(P_i)\ge2^{2k}\ge n\ell$.
Next, suppose that $i:[n]\to\mathcal P$ does form a coloring of $G$. Then there exist at least $\chi(G)$ parts $P_i$ each receiving at least $t$ columns from some column block, and so, since $A$ is $(c,t)$-universal, the contribution from the $A$ columns give us $\operatorname{Cost}(\mathcal P)\ge\chi(G)\cdot2^t=\chi(G)\cdot\ell$.
This will give us the theorem, and all we are left to do is ensure that the various conditions can be met: We define $t$ so that $\ell=2^t$ satisfies that $\ell\le(nm^2)^{1+\epsilon/2}<2\ell$. To meet Condition 3, let $k$ be the smallest integer satisfying that $n\ell\le2^{2k}$. Let $c:=n\ell\cdot\max\{t,k\}$ so that Condition 2 is satisfied. Finally, observe that Condition 1 is satisfied because
\begin{align*}
r+ms&\le\ell^{1+o(1)}\cdot\log c+m\cdot2^{(1+o(1))\cdot k}\cdot\log c\\
&\le\ell^{1+o(1)}+m(n\ell)^{\frac12+o(1)}\cdot\log c\le\ell^{1+\epsilon},
\end{align*}
where the last inequality holds for all large $n,m$.
\end{proof}
We can now prove that $L_3^A$ and $C_3^A$ are hard to approximate, also for Boolean functions.
\begin{theorem}\label{thm:three-round-hardness}
For every constant $\varepsilon>0$, there exists a deterministic quasipolynomial-time algorithm that, on input $x\in\{0,1\}^*$, outputs a communication matrix $M\in\{0,1\}^{N\times N}$ and a number $k\in\mathbb N$, with $k\le N=|x|^{O(1)}$, such that
\begin{enumerate}
\item if $x$ is a YES instance of SAT, then $L_3^B(M)\le k$ and $C_3^B(M)\le\log k$, and
\item if $x$ is a NO instance of SAT, then $L_3^B(M)\ge N^{\frac18-\varepsilon}\cdot k$ and $C_3^B(M)\ge\log k+(\frac18-\varepsilon)\log N$.
\end{enumerate}
\end{theorem}
\begin{proof}
We combine the two reductions of Theorems~9 and~17, which we invoke with the same parameter $\varepsilon$. Fix any input $x$ and let $G$ be an $n$-vertex graph that is produced by the reduction of Theorem~9 on input $x$. Let $M\in\{0,1\}^{a\times b}$ be the communication matrix of $f_G$, and $\ell\in\mathbb N$, be given by the reduction of Theorem~17 on input $G$, and set $N=\max(a,b)=O(n^8)$, $k=n^\varepsilon\cdot\ell^{1+\varepsilon}$. We may assume that $a=b$ without loss of generality, since otherwise we may pad $M$ with all-0 rows or all-0 columns, and this changes the leaf complexity by at most a factor of 2, and the communication complexity by at most 1 bit, so this change makes no difference to the result.
We verify that the inequalities are satisfied for $M$: If $x\in L$, then we have $L_3^B(M)\le\chi(G)\cdot\ell^{1+\epsilon}\le n^\varepsilon\cdot\ell^{1+\epsilon}=k$. If $x\notin L$, then we have
\[L_3^B(M)\ge\chi(G)\cdot\ell\ge n^{1-\epsilon}\cdot\ell\ge N^{\frac18-O(\epsilon)}\cdot k.\]
Similar inequalities hold for $C_3^B(M)$. Since $\varepsilon$ can be arbitrarily small, we may ignore the constant factors.
\end{proof}
